% Generated draft: requires mathematical and source review.
\documentclass[11pt]{article}
\usepackage[margin=25mm]{geometry}
\usepackage{fontspec}
\setmainfont{Latin Modern Roman}
\usepackage{amsmath,amssymb,mathtools,unicode-math}
\setmathfont{Latin Modern Math}
\usepackage{enumitem,needspace,xurl,hyperref}
\setlength{\parindent}{0pt}
\setlength{\parskip}{5pt}
\setlength{\emergencystretch}{3em}
\newcommand{\N}{\mathbb{N}}
\newcommand{\Z}{\mathbb{Z}}
\newcommand{\Q}{\mathbb{Q}}
\newcommand{\R}{\mathbb{R}}
\newcommand{\C}{\mathbb{C}}
\newcommand{\F}{\mathbb{F}}
\newcommand{\A}{\mathbb{A}}
\newcommand{\PP}{\mathbb{P}}
\newcommand{\E}{\mathbb{E}}
\newcommand{\eps}{\varepsilon}
\newcommand{\dd}{\,\mathrm{d}}
\begin{document}
\section*{UnsolvedMath ID: 30002874}
\textbf{Source problem code:} catalogue entry
\subsection{Definitions and mathematical statement}
Let \(Y=(\mathbb{P}^{1}(\C))^{3}\), equipped with coordinatewise complex conjugation \(\sigma\). Under the standard interpretation of a real algebraic surface in this ambient space, \(X\subset Y\) is a nonsingular complex algebraic hypersurface invariant under \(\sigma\), and its real part is
\[
\mathbb{R}X=X^{\sigma}=X\cap(\mathbb{P}^{1}(\R))^{3}.
\]
For a prescribed multidegree \(\mathbf d=(d_1,d_2,d_3)\), let \(\mathcal S_{\mathbf d}\) denote the class of such hypersurfaces of multidegree \(\mathbf d\). Here multidegree means that a defining multihomogeneous equation has degree \(d_i\) in the homogeneous coordinates of the \(i\)-th projective-line factor; equivalently, its divisor line bundle is
\[
\mathcal O_Y(d_1,d_2,d_3)=\bigotimes_{i=1}^{3}\pi_i^{*}\mathcal O_{\mathbb{P}^{1}}(d_i),
\]
where \(\pi_i\) is the projection onto the \(i\)-th factor.

For a specified coefficient field \(K\), define
\[
b_1(\mathbb{R}X;K)=\dim_K H_1(\mathbb{R}X;K).
\]
The extremal question is to determine, for every prescribed multidegree for which \(\mathcal S_{\mathbf d}\) is nonempty, the largest attainable value
\[
M_K(\mathbf d)=\max_{X\in\mathcal S_{\mathbf d}} b_1(\mathbb{R}X;K).
\]
The supplied statement writes \(b_1(\mathbb{R}X)\) without specifying coefficients. Thus this formula makes the coefficient choice explicit but does not resolve the intended choice: rational coefficients and \(\mathbb F_2\) coefficients can give different first Betti numbers for nonorientable real surfaces. The source also does not specify the allowed range of multidegrees or whether irreducibility or connectedness is required. No such additional restrictions are imposed here; their intended role requires review. Coordinatewise conjugation is the standard real structure used in this formulation, but the supplied statement does not explicitly identify the real structure.
\subsection{Short English statement}
For each prescribed multidegree, determine the greatest first Betti number attainable by the real part of a nonsingular real algebraic surface in the product of three complex projective lines. The intended homology coefficients, degree range, and any irreducibility or connectedness requirements need clarification. The supplied registry labels the problem open; this draft does not independently verify that status.
\subsection{Sources}
\begin{itemize}
\item Supplied local UnsolvedMath registry record, permanent ID \(30002874\), title \emph{Maximal Real Betti Numbers of Tridegree Surfaces}, exact source problem code \texttt{catalogue entry}, in the Oberwolfach Reports Open Problems collection. The record lists proposed year 2015 and supplies no named proposer. Its statement status is \texttt{corrected\_verified}, described as preservation of mathematical content with repairs to notation, typography, or omitted definitions; this is supplied registry metadata, not an independent verification here.
\item Supplied report citation: \emph{Tropical Aspects in Geometry, Topology and Physics} (2016), Oberwolfach Reports, DOI 10.4171/owr/2015/23, pp.\ 1255--1316. \url{https://doi.org/10.4171/owr/2015/23}. The record describes these reports as documenting mathematical workshops and their open-problem discussions. Its literature-source label for this link is ``Oberwolfach Report 23.'' The relationship between the citation year 2016 and the supplied proposed year 2015 requires bibliographic review.
\item Supplied literature source: Renaudineau, \emph{Real Algebraic Surfaces with Many Handles}. \url{https://arxiv.org/abs/1511.02261}.
\item Literature assessment checked on \texttt{2026-08-21}, according to the supplied record: status open; constructions for real surfaces in the triple projective line provide counterexamples and asymptotically near-maximal first Betti numbers, while the exact maximum for each prescribed multidegree remains unresolved. This is supplied dated curation derived from the OWR record, not an independently verified literature assessment. The background explicitly requests independent verification of time-sensitive status and cited claims. The assertion about counterexamples does not identify the claim being contradicted.
\end{itemize}
\end{document}
